\documentclass[11pt,a4paper]{article}
\usepackage[T1]{fontenc}
\usepackage[utf8]{inputenc}
\usepackage{lmodern}
\usepackage[margin=27mm]{geometry}
\usepackage{amsmath,amssymb,amsthm,mathtools}
\usepackage{microtype}
\usepackage{enumitem}
\usepackage[authoryear,round]{natbib}
\usepackage{xurl}
\usepackage[colorlinks=true,linkcolor=blue,citecolor=blue,urlcolor=blue]{hyperref}
\newcommand{\workbenchurl}{https://github.com/KokunoYumeto/algebraic-groups-and-representations-workbench}
\newcommand{\recordurl}{https://github.com/KokunoYumeto/mathematics-workbenches}
\newcommand{\Q}{\mathbf Q}
\newcommand{\Z}{\mathbf Z}
\newcommand{\F}{\mathbf F}
\newcommand{\OO}{\mathcal O}
\newcommand{\pp}{\mathfrak p}
\newcommand{\End}{\operatorname{End}}
\newcommand{\Nrd}{\operatorname{Nrd}}
\newcommand{\Cl}{\operatorname{Cl}}
\newcommand{\inv}{\operatorname{inv}}
\newcommand{\ord}{\operatorname{ord}}
\newcommand{\Gal}{\operatorname{Gal}}
\newtheorem{theorem}{Theorem}[section]
\newtheorem{proposition}[theorem]{Proposition}
\newtheorem{lemma}[theorem]{Lemma}
\theoremstyle{definition}
\newtheorem{definition}[theorem]{Definition}
\theoremstyle{remark}
\newtheorem{remark}[theorem]{Remark}
\setlist[enumerate]{leftmargin=*,itemsep=3pt}
\setlength{\parskip}{3pt}
\setlength{\emergencystretch}{2em}
\title{Thakur's descent bound and\protect\\
Grothendieck's finite-field realization question\protect\\[3pt]
\large An abelian threefold, Kummer surfaces, and cubic fourfolds}
\author{}
\date{23 November 2026}
\begin{document}
\maketitle
\vspace{-1.2em}

\begin{abstract}
Thakur's Proposition 2.9 in arXiv:2001.00210v1 proposes that the
class number of an imaginary quadratic endomorphism center bounds the
field needed to define a geometric isogeny class. We give a
counterexample and determine its minimum field exactly. The Weil
$512$-number $12-4\sqrt{-23}$ defines a geometrically simple abelian
threefold. Its geometric isogeny class has a representative over
$\F_{2^e}$ exactly when $9$ divides $e$, although the center has class
number $3$. Honda--Tate theory gives the threefold and its endomorphism
division algebra; a reduced norm, an ideal class, and the local division
algebra exclude smaller fields, including models whose Frobenius is
not initially central. We also give the exact minimum for models with
central Frobenius over a CM center, in terms of an ideal class and a
unit-norm obstruction. This is the full minimum when the geometric
endomorphism algebra is commutative. In the imaginary quadratic case
we recover the sufficient class-number times Brauer-index bound from
Thakur's earlier construction. A finite polynomial procedure determines the minimum
exponent for any given abelian variety. A companion section gives
explicit Chow retractions for Kummer surfaces and separated cubic
fourfolds, and constructs abelian targets for ordinary K3 surfaces
and cubic fourfolds. For the latter, the missing comparison
correspondence is characterized by a specific mixed Tate condition.
These results concern concrete aspects of Grothendieck's
finite-field notes; they do not decide his broader existence
conjecture. GPT-6 Astra (ChatGPT Pro mode) produced the text
and mathematics and performed the same-model checks; no human has
checked them.
\end{abstract}

\section{The descent question and the result}

The field of definition of a geometric isogeny class records more than
the ideal class of a Frobenius number. Its local valuations must also
give the correct endomorphism algebra and dimension. We make this
distinction explicit in a three-dimensional example.

For an abelian variety $A$ over a finite field of characteristic $p$,
write $A_{\overline{\F}_p}$ for its base change to an algebraic closure.
We say its \emph{geometric isogeny class has a model over $\F_{p^e}$}
if some abelian variety $B/\F_{p^e}$ is isogenous to
$A_{\overline{\F}_p}$ after base change to $\overline{\F}_p$.
The minimum such positive integer $e$ is the \emph{minimum exponent}.
No polarization or prescribed finite-field isogeny is required.

\begin{theorem}\label{thm:main}
Put $K=\Q(s)$, where $s^2=-23$, and $\gamma=12-4s$.
There is a simple abelian variety $A/\F_{512}$ with Frobenius
$\gamma$, unique up to $\F_{512}$-isogeny. It has the following
properties.
\begin{enumerate}
\item It is a geometrically simple threefold, and its Frobenius
characteristic polynomial is
\[
                 (T^2-24T+512)^3.
\]
\item Its geometric rational endomorphism algebra is a central
division algebra $D$ of degree $3$ over $K$. Its invariants are
$2/3$ and $1/3$ at the two primes above $2$, and zero everywhere else.
\item The class number of $K$ is $3$.
\item For every positive integer $e$, there is an abelian variety
$B/\F_{2^e}$ geometrically isogenous to $A$ if and only if $9\mid e$.
\end{enumerate}
In particular the minimum exponent is $9$, and no such model exists
over $\F_8$.
\end{theorem}

In the terminology of \citet[Definition 2.1, p.~3]{Thakur2020}, this
is an abelian variety of type $\mathrm{IV}(1,3)$.
His Definition 2.3 and Proposition 2.9 (p.~6) assert a model up to
isogeny over $\F_{p^{h_K}}$ for type $\mathrm{IV}(1,d)$, where
$h_K$ is the class number of the center. Theorem~\ref{thm:main}
contradicts even the weaker geometric version of that assertion.
The arXiv record consulted on 23 November 2026 lists only version 1,
dated 1 January 2020. The correction here refers to that statement.

The relevant positive construction is already in
\citet[Proposition 2.4, pp.~4--5]{Thakur2020}: principalize a prime
above $p$, and combine a generator and its conjugate with exponents
that preserve the desired slopes. Proposition~\ref{prop:central}
below formulates its geometric consequence and computes the least
exponent when the Frobenius is required to be central. The example
shows that the resulting upper bound $h_Kd$ can be attained.

For clarity, the general central formula has the following full
scope. Let a geometrically simple class have geometric division
algebra $D$ of degree $d$ over a CM field $K$, with stable slopes
$\lambda_v$ at $v\mid p$. Put
$a=\operatorname{lcm}_{v\mid p}\operatorname{den}(e_v\lambda_v)$,
where $e_v$ is the ramification index, and set
$I=\prod_{v\mid p}\pp_v^{ae_v\lambda_v}$. If $t$ is the order
of $[I]$, choose $(\eta)=I^t$ and let $s$ be the order of
$\eta\bar\eta/p^{at}$ in the group of totally positive units
of $K^+$ modulo norms of units of $K$. Then $s=1$ or $2$, and
the central-Frobenius exponents are exactly the multiples of $ats$.
This is the unrestricted minimum when $D=K$.
Proposition~\ref{prop:central-cm} proves the formula and the
sufficient bound $2h_Kd\operatorname{lcm}_{v\mid p}f_v$, with
$f_v$ the residue degrees. These are the familiar ideal and unit
groups of polarized ideal-class theory
\citep[Lemma 5.5 and \S5.3]{BrooksEtAl2017}; the bound is not
claimed sharp.

Section~\ref{sec:algorithm} gives an explicit finite test for the
minimum exponent when a Frobenius polynomial is given. This is a
transparent variant of established computational methods, not a claim
of a new algorithm: Weil-polynomial enumeration and the
root-of-unity test for geometric isogeny are described by
\citet[\S\S3.1 and 3.4, pp.~12--15]{DKRV2020}.

\paragraph{Grothendieck's question and the scope.}
In cote 15, scan 100 (manuscript footer 99), Grothendieck asks whether
the arithmetic conditions he lists on Frobenius eigenvalues are
sufficient for geometric realization. Scans 102--104 (footers
101--103) discuss passage to powers under finite-field extension;
scan 106 (footer 105) states a weaker geometric existence conjecture
\citep{GrothendieckND}. These locators are checked against the
transcription, whose uncertain readings are not needed here. The
present result concerns the weight-one setting: abelian varieties up
to isogeny, identified with effective weight-one motives in
\citet[\S1, p.~95]{Tate1971}. It refutes the later descent bound,
not Grothendieck's existence conjecture. We do not prove a closed
minimum formula allowing noncentral Frobenius over arbitrary CM
centers, or a sharp bound there depending only on class numbers and
Brauer indices.

Section~\ref{sec:companion} addresses a complementary part of the
realization problem. It distinguishes explicit Chow retractions,
the existence of an abelian target with the right Frobenius
realization, and the mixed algebraic correspondence needed to
identify a given variety with that target. The Kummer and
separated-cubic constructions use classical geometric arguments;
the general comparison criterion keeps its unproved input visible.

Here is the full companion statement. A Chow retraction consists
of rational algebraic correspondences whose composition is the
diagonal modulo rational equivalence. A homological retraction
asserts the corresponding identity on cohomology. Write
$\mathsf h(X)$ for the motive of $X$ and
$\mathbb L=\Q(-1)$ for the Lefschetz motive. The Tate conjecture
below means surjectivity of the rational $\ell$-adic cycle class
map onto Frobenius invariants, for every $\ell\ne p$ unless one
prime is explicitly specified.

\begin{theorem}\label{thm:companion-main}
The following statements hold over finite fields.
\begin{enumerate}
\item If $A$ is an abelian surface in characteristic different from
two, its Kummer surface $K$ has a Chow retraction from
$\mathsf h(A)\oplus\mathsf h(A[2])\otimes\mathbb L$.
If $A=E_1\times E_2$, all powers of $K$ satisfy Tate in every
codimension after every finite extension. Its formal Brauer height
is one when both $E_i$ are ordinary, and two when exactly one is
ordinary and the other is supersingular.
\item In characteristic different from three, let $f$ and $g$ be
ternary cubic forms whose plane curves $E_f,E_g$ are smooth.
The fourfold $X:(f+g=0)\subset\mathbf P^5$ is smooth, and after
a finite extension its Chow motive is a summand of
\[
 (\mathbf1\oplus\mathbb L^{\oplus7}\oplus\mathbb L^2)^{\otimes2}
       \oplus\mathsf h(E_f\times E_g)\otimes\mathbb L.
\]
It is a Chow summand of abelian, Artin, and Lefschetz motives over
the original field. All its powers satisfy Tate in every
codimension after every finite extension.
\item Let $X$ be an ordinary K3 surface or an ordinary smooth
cubic fourfold, in any characteristic. After a finite extension,
there is an abelian variety $B$ whose powers satisfy Tate in every
codimension after every finite extension, and a rational
homological projector $e$ on a finite direct sum
$T=\bigoplus_\nu\mathsf h(Y_\nu)$ of powers of $B$, such that
$M=(T,e)$ has the same Frobenius characteristic polynomial as
$X$ in every degree. There are rational algebraic correspondences
$\mathsf h(X)\rightleftarrows M$ inverse in every
$\ell$-adic realization if and only if, for one $\ell_0\ne p$,
\[
 \bigoplus_\nu\operatorname{CH}^{\dim X}(X\times Y_\nu)_\Q
            \otimes\Q_{\ell_0}
 \longrightarrow
 \operatorname{Hom}_{F}^{\mathrm{gr}}
       (H^*_{\ell_0}(X),H^*_{\ell_0}(M)),\qquad z\longmapsto e_*z_*,
\]
is surjective. Equivalently, Tate in codimension $\dim X$ holds
at one prime for each $X\times Y_\nu$.
\end{enumerate}
The equivalence in (3) also holds for any finite-height K3 surface
or any smooth cubic fourfold whenever such a target is given in a
tensor category of abelian, Artin, and Lefschetz motives on whose
generating products Tate holds in every codimension and after every
finite extension. Existence of that target outside the ordinary
case is an additional hypothesis here.
\end{theorem}

The first two parts supply algebraic correspondences; the third
constructs a target and identifies the precise remaining condition.
The cubic construction is the degree-three case of
\citet[Remark 1.10, p.~102]{ShiodaKatsura1979}. The all-powers
Tate inputs for the abelian targets are the two examples in
\citet[Appendix A.7, pp.~158--159]{Milne2001}. The newer Tate
results for K3 surfaces and cubic fourfolds enter only where
specified in Section~\ref{sec:companion}. In particular, the
self-product theorems of \citet[Theorem 1.5]{IIK2021} and
\citet[Corollary 5.2]{Zarhin2004} do not supply the mixed
correspondences required in (3).

\section{Proof overview}

The norm of $\gamma$ gives its complex absolute values. The two
valuations above $2$ give slopes $2/3$ and $1/3$, so Honda--Tate
theory attaches a threefold with division algebra of degree three.
Those slopes persist after every extension and keep the Frobenius
field quadratic; this proves geometric simplicity and identifies all
geometric endomorphisms.

To exclude smaller fields, it is not enough to look for a root of
$\gamma$ in $K$: a model could in principle have a noncentral
Frobenius in $D$. We therefore transport any proposed Frobenius into
$D$. Its reduced norm is central, and the nonprincipal prime class
above $2$ forces the exponent to be divisible by three. The remaining
case would give an element whose cube is a signed power of $\gamma$.
Every such power has a cube root in $\Q_2$, whereas a degree-three
division algebra has no quadratic subfield. This forces the proposed
Frobenius itself to be central. The same ideal class then forces a
second factor of three. Controlling roots of unity in possible
subfields of $D$ makes the argument cover all twists.

For a general CM center, the integral slope denominator determines
an ideal. Its class must first become principal; its generator must
then have the correct norm to the real subfield. The remaining
unit class has order at most two. Comparison with the original
Frobenius by a root of unity shows that this arithmetic construction
preserves the geometric isogeny class.

For the companion results, blowups and generically finite maps give
Chow retractions. In the ordinary case, the unique slope-zero root
forces an irreducible, power-stable CM field. Honda--Tate theory
then produces an abelian variety with division-algebra index at
most two. Squaring that Brauer class and using an algebraic
idempotent removes the cohomological multiplicity. A mixed Tate
condition supplies a forward comparison correspondence; duality
on the target supplies its inverse. Rational trace identities make
one-prime invertibility a certificate valid at all primes.

\section{The arithmetic and the Honda--Tate variety}

This section computes the Frobenius data and verifies that the
endomorphism algebra does not enlarge over the algebraic closure.
Let
\[
 \omega=\frac{1+s}{2},\qquad
 \OO_K=\Z[\omega],\qquad \omega^2-\omega+6=0,
 \qquad \pp=(2,\omega),\quad
 \bar\pp=(2,\omega-1).
\]
The discriminant of $K$ is $-23$. Reduction of the displayed
polynomial modulo $2$ gives two distinct roots, so
$(2)=\pp\bar\pp$, and both completions are $\Q_2$. Since
\[
 \gamma=8(2-\omega),\qquad
 (2-\omega)(1+\omega)=8,
\]
and $2-\omega$ vanishes at $\pp$ but not at $\bar\pp$, we obtain
\begin{equation}\label{eq:ideal}
 (2-\omega)=\pp^3,
 \qquad (\gamma)=\pp^6\bar\pp^3,
 \qquad \gamma\bar\gamma=512,
 \qquad \gamma+\bar\gamma=24.
\end{equation}

\begin{lemma}\label{lem:class}
The ideal class of $\pp$ has order three, and $h_K=3$.
\end{lemma}
\begin{proof}
An element of $\OO_K$ is $(a+bs)/2$ with integers $a,b$ of the
same parity, and its norm is $(a^2+23b^2)/4$. Norm two would
require $a^2+23b^2=8$, which is impossible. Thus $\pp$ is not
principal, whereas \eqref{eq:ideal} makes $\pp^3$ principal.

The imaginary quadratic Minkowski bound is
$2\sqrt{23}/\pi<4$ \citep[Theorem 4.3, p.~70]{Milne2020}.
Every ideal class therefore has a representative of norm $1$, $2$,
or $3$. At $3$ put $\mathfrak q=(3,\omega)$ and
$\bar{\mathfrak q}=(3,\omega-1)$. The polynomial of $\omega$
again has distinct roots modulo $3$, and $N_{K/\Q}(\omega)=6$
gives $(\omega)=\pp\mathfrak q$. Thus
$[\mathfrak q]=[\pp]^{-1}$; conjugation gives the other class.
Consequently the only classes are $1,[\pp],[\pp]^2$.
\end{proof}

We use the Honda--Tate theorem in the following precise form
\citep[Theorem 1(i)--(ii) and Remark 2, p.~96]{Tate1971}.
An integral Weil $q$-number $\pi$ determines a simple isogeny class
over $\F_q$. If $E=\Q(\pi)$ and $D_\pi$ is its endomorphism
division algebra, then for $v\mid p$,
\begin{equation}\label{eq:HT}
 \inv_v(D_\pi)
   =[E_v:\Q_p]\frac{v(\pi)}{v(q)}\pmod{\Z},
 \qquad
 2\dim A_\pi=[E:\Q]\sqrt{[D_\pi:E]}.
\end{equation}
The finite invariants away from $p$ vanish. At a real place the
invariant is $1/2$; there are no real places in our application.
The degree $\sqrt{[D_\pi:E]}$ is the least common denominator of
the local invariants, and the characteristic polynomial is the
minimal polynomial of $\pi$ to that power. We use Tate's positive
sign convention for endomorphisms of abelian varieties.

By \eqref{eq:ideal}, $\gamma$ is an integral Weil $512$-number
with minimal polynomial $T^2-24T+512$. Formula \eqref{eq:HT}
gives the two invariants $6/9=2/3$ and $3/9=1/3$. The division
algebra has degree three, hence the abelian variety has dimension
three and the polynomial asserted in Theorem~\ref{thm:main}.

For every $n\geq1$, the valuations of $\gamma^n$ at the two
primes are $6n$ and $3n$. A rational number has equal valuations
there, so $\Q(\gamma^n)=K$. Relative to $2^{9n}$ its invariants
remain $2/3$ and $1/3$. The corresponding simple isogeny class again
has dimension three. The Frobenius polynomial of
$A_{\F_{2^{9n}}}$ is the cube of the minimal polynomial of
$\gamma^n$, so the Honda--Tate classification identifies it with
this simple class. Thus $A$ remains simple over every finite
extension and is geometrically simple.

The endomorphism algebras over these finite extensions all have
$\Q$-dimension $2\cdot3^2=18$. Their natural inclusions from
$\End^0_{\F_{512}}(A)$ are therefore equalities. Every geometric
endomorphism is defined over some finite extension, so
\[
                     \End^0_{\overline{\F}_2}(A)=D.
\]
This proves the first three assertions of the theorem.

\section{All possible fields of definition}

We now include possible noncentral Frobenius elements and exclude
them where they could lower the minimum exponent.

\begin{lemma}\label{lem:torsion}
If $L$ is a subfield of $D$ containing $K$, then $[L:K]$ is
$1$ or $3$, and the only roots of unity in $L$ are $1$ and $-1$.
\end{lemma}
\begin{proof}
Every subfield of a central division algebra is contained in a
maximal subfield. A maximal subfield of an algebra of degree three
has degree three over the center; see
\citet[Lemma 11.7.4, Tag 074W]{Stacks}. Hence $[L:\Q]$ is $2$
or $6$. If a primitive $r$th root of unity belongs to $L$, then
$\varphi(r)$ divides $6$. The possible orders are
\[
                        1,2,3,4,6,7,9,14,18.
\]
Orders divisible by $3$ or $4$ would give the quadratic subfield
$\Q(\sqrt{-3})$ or $\Q(i)$. Its compositum with $K$ would have
degree four over $\Q$, impossible inside a field of degree two or
six. For order $7$ or $14$, the field would be $\Q(\zeta_7)$.
Its unique quadratic subfield is $\Q(\sqrt{-7})$: indeed
$u=\zeta_7+\zeta_7^2+\zeta_7^4$ satisfies $u^2+u+2=0$,
and the cyclic group $(\Z/7\Z)^\times$ has one subgroup of
index two. This also excludes $K\subset\Q(\zeta_7)$.
Only orders $1$ and $2$ remain.
\end{proof}

\begin{lemma}\label{lem:local}
Let $D_2$ be a central division algebra of degree three over
$\Q_2$. If $x\in D_2$ and $x^3=c^3$ for some nonzero
$c\in\Q_2$, then $x=c$.
\end{lemma}
\begin{proof}
Because $c$ is central,
\[
                 (x-c)(x^2+cx+c^2)=0.
\]
If $x\ne c$, the absence of zero divisors gives
$(x/c)^2+x/c+1=0$. The polynomial $T^2+T+1$ is irreducible
over $\Q_2$, since its reduction is irreducible over $\F_2$.
It would define a quadratic subfield of $D_2$, which is impossible
by the maximal-subfield degree formula just used.
\end{proof}

\begin{proof}[Proof of Theorem~\ref{thm:main}, assertion (4)]
Suppose $B/\F_{2^e}$ is geometrically isogenous to $A$. Through
such an isogeny, its Frobenius gives an element $\beta\in D$.
The isogeny is defined over some finite field. Enlarging this field,
we may take it to be $\F_{2^{9em}}$ for a positive integer $m$.
Compatibility with Frobenius gives
\begin{equation}\label{eq:common}
                         \beta^{9m}=\gamma^{em}.
\end{equation}
Put $a=\Nrd_{D/K}(\beta)\in K$. Reduced norm is
multiplicative and sends a central element $z$ to $z^3$. Thus
\[
              a^{9m}=\gamma^{3em},\qquad
              (a^3\gamma^{-e})^{3m}=1.
\]
The roots of unity in $K$ are $\{\pm1\}$, so
$a^3=\pm\gamma^e$. Taking fractional ideals in $K$ and using
\eqref{eq:ideal} yields
\[
             (a)=\pp^{2e}\bar\pp^e=(2)^e\pp^e.
\]
Lemma~\ref{lem:class} now implies $3\mid e$.

Write $e=3k$. The subalgebra $L=K(\beta)$ is a field:
$\beta$ is algebraic, and a finite commutative subalgebra of a
division algebra has no zero divisors. Equation \eqref{eq:common}
implies
\[
                       (\beta^3\gamma^{-k})^{3m}=1.
\]
Lemma~\ref{lem:torsion} gives $\beta^3=\pm\gamma^k$.
At either place $v\mid2$, the valuation of $\pm\gamma^k$ in
$K_v=\Q_2$ is divisible by three. Every $2$-adic unit is a cube:
Hensel's lemma applies to the simple root $1$ modulo $2$ of
$T^3-u$ for $u\in\Z_2^\times$. Therefore
$\pm\gamma^k=c^3$ for some $c\in\Q_2^\times$.

The completion $D_v$ is a division algebra of degree three because
its invariant has denominator three. Lemma~\ref{lem:local} shows
that the image of $\beta$ in $D_v$ is central. Consequently
$\beta$ is central in $D$: for each $z\in D$, the commutator
$\beta z-z\beta$ becomes zero in $D_v$, and the scalar-extension
map $D\longrightarrow D_v$ is injective. Hence $\beta\in K$.
Finally, $\beta^3=\pm\gamma^k$ gives
\[
                    (\beta)=(2)^k\pp^k.
\]
The order-three ideal class forces $3\mid k$, so $9\mid e$.
Conversely, when $9\mid e$, base change of $A/\F_{2^9}$ to
$\F_{2^e}$ supplies the required model.
\end{proof}

\section{The bound with the Brauer index}

The example fits a general construction in which the ideal class
and the Newton slope enter separately.

\begin{proposition}\label{prop:central}
Let $X$ be a geometrically simple abelian variety over a finite
field of characteristic $p$. Suppose its geometric endomorphism
algebra is a central division algebra $D$ of degree $d>1$ over an
imaginary quadratic field $K$, and $p$ splits as
$\pp\bar\pp$ in $K$. Write
\[
 \inv_{\pp}(D)=\frac jd,\qquad
 0<j<d,\quad \gcd(j,d)=1,
 \qquad t=\ord_{\Cl(K)}[\pp].
\]
Then $2j\ne d$. Among models in the geometric isogeny class of
$X$ whose Frobenius lies in the center of the geometric
endomorphism algebra, the possible exponents are exactly the
positive multiples of
\begin{equation}\label{eq:central-min}
                       n_0=\frac{dt}{\gcd(t,2j-d)}.
\end{equation}
In particular the geometric isogeny class has a model over
$\F_{p^{td}}$ and over $\F_{p^{h_Kd}}$.
\end{proposition}

\begin{proof}
Choose $N$ so that $X$ and a $\Q$-basis of all its geometric
rational endomorphisms are defined over $\F_{p^N}$. Its Frobenius
$\pi$ commutes with these endomorphisms, hence lies in $K$.
Since $X$ is simple, Honda--Tate gives $\Q(\pi)=K$ and
$\dim X=d$. The two normalized valuations are $j/d$ and
$(d-j)/d$: they lie in $[0,1]$ and represent the indicated
nonzero local invariants. Thus $d\mid N$.

If $2j=d$, the two valuations would be equal. Then
$\pi^2/p^N$ would be a unit in the imaginary quadratic field and
therefore a root of unity. Some power of $\pi$ would be rational.
Over that extension, Honda--Tate would give a rational center,
contradicting the geometric center $K$. Thus $2j\ne d$.

For any central Frobenius model over $\F_{p^n}$, every geometric
endomorphism commutes with Frobenius and hence is defined over
that field. Its finite-field endomorphism algebra is therefore
$D$. The same invariant calculation forces $n=dr$ and
\begin{equation}\label{eq:principal}
       (\pi_n)=\pp^{jr}\bar\pp^{(d-j)r}.
\end{equation}
The ideal on the right is principal precisely when
$t\mid(2j-d)r$, since $[\bar\pp]=[\pp]^{-1}$.
This is equivalent to $n_0\mid n$.

Conversely, suppose $n=dr$ and $t\mid(2j-d)r$. Choose a
generator $\rho$ of the integral ideal on the right of
\eqref{eq:principal}. Its norm is $p^n$, so it is an integral Weil
$p^n$-number. Unequal valuations imply
$\Q(\rho^a)=K$ for every $a>0$. Its invariants have denominator
$d$, so Honda--Tate gives a geometrically simple variety $X_\rho$
of dimension $d$, with geometric endomorphism algebra $D$ and
central Frobenius $\rho$.

The elements $\rho^N$ and $\pi^n$ have equal valuations at all
finite primes of $K$. Their quotient is a unit in $K$, hence has
finite order, say $b$. Thus $\rho^{Nb}=\pi^{nb}$.
The two varieties have the same Frobenius polynomial after
extension to $\F_{p^{Nnb}}$, and are isogenous there by
Honda--Tate. This proves sufficiency. Since $n_0\mid dt$ and
$t\mid h_K$, the two stated fields are available.
\end{proof}

The bound $td$ comes directly from the generator construction of
\citet[Proposition 2.4, pp.~4--5]{Thakur2020}. Explicitly, if
$(\eta)=\pp^t$, then
$\eta^j\bar\eta^{d-j}$ has the required slopes over
$\F_{p^{td}}$. Formula \eqref{eq:central-min} also allows
cancellation in the ideal class. In Theorem~\ref{thm:main},
$d=t=h_K=3$ and $j=2$, so $n_0=9$. Section 4 proves that
allowing noncentral Frobenius cannot lower this value. Hence the
upper bound $h_Kd$ is sharp in the stated imaginary quadratic
setting.

The precise failed inference in the class-number-only argument is
visible in the smallest calculation. Here
$\eta=2-\omega=(3-s)/2$ generates $\pp^3$ and is a Weil
$8$-number. Its two valuations are $3$ and $0$, so its slopes are
$1$ and $0$. Its Honda--Tate division algebra has degree one and
its abelian variety is an ordinary elliptic curve, with polynomial
$T^2-3T+8$. It cannot be assigned the degree-three algebra of
$A$. Principalizing a prime ideal alone loses the slope information.

\subsection{Central Frobenius over a general CM center}

The imaginary quadratic calculation extends to CM centers by adding
a unit-norm obstruction. The relevant ideal and unit groups are the
ones used in polarized ideal-class theory; compare
\citet[Lemma 5.5 and \S5.3]{BrooksEtAl2017}. We give the argument
directly for the descent question.

\begin{proposition}\label{prop:central-cm}
Let $A/\overline{\F}_p$ be a geometrically simple abelian variety, with geometric
endomorphism division algebra $D$ of degree $d$ over a CM field
$K$. Choose a model over $\F_{p^N}$ on which all geometric
endomorphisms are defined, with Frobenius $\pi_0\in K$. For
$v\mid p$, write $e_v=\ord_v(p)$, let $f_v$ be the residue degree,
and put $\lambda_v=\ord_v(\pi_0)/(Ne_v)$. Define
\[
 a=\operatorname{lcm}_{v\mid p}\operatorname{den}(e_v\lambda_v),
 \qquad
 I=\prod_{v\mid p}\pp_v^{a e_v\lambda_v},
 \qquad t=\ord_{\Cl(K)}[I],
\]
where $\operatorname{den}$ is the positive reduced denominator.
Choose $(\eta)=I^t$, and let
\[
 u=\frac{\eta\bar\eta}{p^{at}}\in\OO_{K^+}^{\times,+},
 \qquad
 s=\ord\left([u]\in
 \OO_{K^+}^{\times,+}/N_{K/K^+}(\OO_K^\times)\right).
\]
Here $K^+$ is the maximal totally real subfield and the superscript
$+$ means totally positive. Then $s\in\{1,2\}$, independently
of the choice of $\eta$. The possible central-Frobenius model
exponents are exactly the positive multiples of
\begin{equation}\label{eq:cm-central-min}
                         n_{\mathrm{cent}}=ats.
\end{equation}
Moreover
\[
 a\mid d\operatorname{lcm}_{v\mid p}f_v,
 \qquad n_{\mathrm{cent}}\mid 2h_Ka
               \mid 2h_Kd\operatorname{lcm}_{v\mid p}f_v.
\]
If $D=K$, formula \eqref{eq:cm-central-min} is the minimum among
all models in the geometric isogeny class.
\end{proposition}

\begin{proof}
The Weil relation gives $\lambda_v+\lambda_{\bar v}=1$ and
$e_v=e_{\bar v}$, whence $I\bar I=(p)^a$. Thus the displayed
$u$ is a totally positive unit of $K^+$. Replacing $\eta$ by
$\eta\varepsilon$ multiplies $u$ by the norm of a unit. For
every totally positive unit $z$ of $K^+$ one has
$z^2=N_{K/K^+}(z)$. Consequently the unit quotient has exponent
at most two, proving the assertions about $s$.

Let $\beta\in K$ be the Frobenius of a compatible central model
over $\F_{p^n}$, transported into $D$ by a geometric isogeny.
Over a common finite field, powers of $\beta$ and $\pi_0$ agree.
Their slopes are therefore equal under this identification of
$K$. The integer valuations $\ord_v(\beta)=ne_v\lambda_v$
force $a\mid n$. Write $n=am$. Then $(\beta)=I^m$, so
$t\mid m$. Writing $m=tr$, there is a unit $\varepsilon$ with
$\beta=\eta^r\varepsilon$. The Weil condition is precisely
\[
       \beta\bar\beta=p^{atr}
       \quad\Longleftrightarrow\quad
       u^rN_{K/K^+}(\varepsilon)=1.
\]
It is possible exactly when $s\mid r$. This proves necessity.

Conversely choose a unit $\varepsilon$ with
$N_{K/K^+}(\varepsilon)=u^{-s}$ and set
$\beta=\eta^s\varepsilon$, $n=ats$. The ideal of $\beta$ is
integral and $\beta\bar\beta=p^n$, so $\beta$ is a Weil
$p^n$-number. The quotient $\beta^N/\pi_0^n$ has valuation zero
at every finite prime and absolute value one at every complex
embedding. It is a root of unity, say of order $b$. Hence
$\beta^{Nb}=\pi_0^{nb}$.

Every positive power of $\pi_0$ generates $K$: all endomorphisms
remain defined after extension, so Honda--Tate identifies the
center there with $\Q(\pi_0^r)=K$. The equality of powers now
also gives $\Q(\beta^r)=K$ for every $r>0$. The slopes of
$\beta$ have the same local invariants as those of $\pi_0$.
Honda--Tate therefore gives the same division-algebra degree and
dimension, and the equality of powers gives a geometric isogeny.
This constructs the central model at $n=ats$; base extension gives
all its positive multiples.

Finally $[K_v:\Q_p]=e_vf_v$, and the local invariant of $D$ is
$e_vf_v\lambda_v$ modulo $\Z$. Its denominator divides $d$.
Thus the denominator of $e_v\lambda_v$ divides $df_v$, proving
the divisibility for $a$. Now $t\mid h_K$ and $s\mid2$ give
the two remaining divisibilities. If $D=K$, every transported
Frobenius is central, so there are no additional models to consider.
\end{proof}

For an imaginary quadratic center with split $p$ as in
Proposition~\ref{prop:central}, one has $a=d$,
$I=\pp^j\bar\pp^{d-j}$, and $K^+=\Q$. The positive unit
quotient is trivial, while
\[
 \ord[I]=\frac{\ord[\pp]}{\gcd(\ord[\pp],2j-d)}.
\]
Thus the general formula recovers \eqref{eq:central-min}.
For a higher-degree center, the residue degrees and the unit-norm
class are explicit additional data. The bound just proved is not
asserted to be sharp. When $D$ is noncommutative, a noncentral
Frobenius may give a smaller exponent; the formula does not exclude
that possibility.


\section{Computing the minimum from Frobenius}\label{sec:algorithm}

The absence of a general closed formula does not prevent an exact
finite computation for a given variety. We spell out such a procedure
to distinguish these two questions. It is not intended to be efficient.

For a monic polynomial $P(T)=\prod_i(T-\alpha_i)$ and an integer
$a>0$, put
\[
 P^{[a]}(T)=\prod_i(T-\alpha_i^a)
           =\operatorname{Res}_U(P(U),T-U^a).
\]
This is the characteristic polynomial after extension of the base
field by degree $a$. This operation and geometric-isogeny tests are
standard; see \citet[\S3.4, pp.~14--15]{DKRV2020}.

\begin{proposition}\label{prop:finite-test}
Let $A/\F_{p^n}$ and $B/\F_{p^e}$ have the same positive dimension
$g$, with Frobenius polynomials $P$ and $Q$. Define
\[
 M_g=\operatorname{lcm}\{r\ge1:\varphi(r)\le4g^2\},
 \qquad l=\operatorname{lcm}(n,e).
\]
Then $A$ and $B$ are geometrically isogenous if and only if
\begin{equation}\label{eq:finite-test}
       P^{[(l/n)M_g]}=Q^{[(l/e)M_g]}.
\end{equation}
Consequently the minimum exponent of the geometric isogeny class
of a given $A/\F_{p^n}$ can be computed by finitely many exact
polynomial and number-field operations.
\end{proposition}

\begin{proof}
First the definition of $M_g$ is effective and finite. For every
$r=\prod_\ell\ell^{a_\ell}$ one has
\[
 \frac{\varphi(r)^2}{r}
   =\prod_{\ell\mid r}\ell^{a_\ell-2}(\ell-1)^2
   \ge \frac12.
\]
The only factor smaller than one can be the factor $1/2$ at $\ell=2$,
$a_\ell=1$. Thus $\varphi(r)\le4g^2$ implies $r\le32g^4$,
giving a finite list for the least common multiple.

Write the roots of $P$ and $Q$, with multiplicity, as $\alpha_i$ and
$\beta_i$. A geometric isogeny is defined over some finite common
extension $\F_{p^{lk}}$. The finite-field isogeny theorem then
matches the roots so that
\[
       (\alpha_i^{l/n})^k=(\beta_i^{l/e})^k.
\]
Their ratio is a root of unity in $\Q(\alpha_i,\beta_i)$, a
number field of degree at most $(2g)^2$. If its order is $r_i$,
then $\varphi(r_i)\le4g^2$, so $r_i\mid M_g$. The same
matching therefore gives \eqref{eq:finite-test}. Conversely that
equality gives an isogeny over $\F_{p^{lM_g}}$ by the
finite-field isogeny theorem. Multiplicities are preserved throughout.

Here is a finite procedure. For $e=1,2,\ldots,n$, put $q=p^e$
and enumerate all monic integral polynomials
\[
 Q(T)=T^{2g}+c_1T^{2g-1}+\cdots+c_{2g},
 \qquad |c_i|\le {2g\choose i}q^{i/2}.
\]
The coefficient bounds follow from the elementary symmetric
functions of roots of absolute value $\sqrt q$. Test exactly which
polynomials have all their complex roots of that absolute value;
algebraic root isolation makes this an exact finite test. Factor each
such polynomial over $\Q$ as $\prod_j f_j^{b_j}$. For a root
$\rho_j$ of $f_j$, compute the Honda--Tate division-algebra degree
$d_j$ from its local invariants, including the real places when
present. Retain precisely the polynomials for which $d_j\mid b_j$
for every $j$. Honda--Tate theory says exactly that these are the
Frobenius polynomials of dimension-$g$ abelian varieties over
$\F_q$: the simple polynomial associated with $\rho_j$ is
$f_j^{d_j}$, and products correspond to products of simple factors
\citep[Theorem 1 and Remark 2, p.~96]{Tate1971}.

Apply \eqref{eq:finite-test} to each retained $Q$. The first
exponent that succeeds is the minimum. The original polynomial
$P$ ensures success by $e=n$, so the procedure terminates. All
factorizations, local valuations, field degrees, real-place tests,
and polynomial equalities used here are effective exact operations.
\end{proof}

For actual computations, the coefficient enumeration and the
cyclotomic resultant tests in \citet[\S\S3.1 and 3.4]{DKRV2020}
avoid much of this deliberately coarse search. Proposition
\ref{prop:finite-test} supplies a uniform finite bound on the
extension needed for each comparison; it gives no useful general
estimate for the smallest field in terms of $h_K$ and the Brauer
index alone.

\section{Companion results on algebraic correspondences}\label{sec:companion}

We now separate three different assertions suggested by the finite-field
realization problem. Some varieties admit explicit retractions in
rational Chow motives. For other varieties, an abelian motive with the
right Frobenius realization can be constructed, but an algebraic
correspondence from the original variety remains to be found. Finally,
for a specified target, that remaining assertion is equivalent to a
precise mixed Tate-conjecture statement. None of these distinctions is
removed by the descent calculation above.

Write $\mathsf h(X)$ for the motive of a smooth projective variety and
$\mathbb L=\Q(-1)$ for the Lefschetz motive, with Frobenius eigenvalue
$q$ over $\F_q$. All correspondences have rational coefficients.
We use the convention
\[
 \operatorname{Hom}(\mathsf h(X),\mathsf h(Y))
       =\operatorname{CH}^{\dim X}(X\times Y)_{\Q}
\]
for untwisted motives; the induced map preserves cohomological degree.
The first two constructions below take place in Chow motives, so their
identities hold modulo rational equivalence. The subsequent target and
comparison statements take place modulo homological equivalence and
assert identities in every $\ell$-adic realization, $\ell\ne p$.
They do not assert Chow isomorphisms.

\subsection{Retractions and finite-field descent}

Two elementary operations turn geometric maps over an extension
field into motivic summands over the original field.

\begin{lemma}\label{lem:retractions}
Let $X$ and $V$ be smooth projective varieties of the same dimension.
If $\phi:V\to X$ is a proper generically finite morphism of degree $a$,
then $\mathsf h(X)$ is a Chow summand of $\mathsf h(V)$, with inclusion
$\phi^*$ and retraction $a^{-1}\phi_*$. If $Z\subset Y$ is smooth of
codimension two, then
\[
 \mathsf h(\operatorname{Bl}_Z Y)
       \simeq \mathsf h(Y)\oplus\mathsf h(Z)\otimes\mathbb L.
\]
\end{lemma}

\begin{proof}
The projection formula gives $\phi_*\phi^*=a$ on Chow groups after
product with every smooth projective variety. Equivalently, composition
of the two graph correspondences is $a\Delta_X$, which proves the
first assertion. For the second, let $b:\widetilde Y\to Y$ be the
blowup and $j:E\hookrightarrow\widetilde Y$, $\rho:E\to Z$ its
exceptional divisor and projection. The maps are
\[
 (b^*,j_*\rho^*),\qquad (b_*,-\rho_*j^*).
\]
The identities $b_*b^*=1$, $b_*j_*\rho^*=0$,
$\rho_*j^*b^*=0$, and
$-\rho_*j^*j_*\rho^*=1$ follow from the projection formula and
$\mathcal O_E(E)=\mathcal O_E(-1)$. The blowup formula says that
these maps exhaust the Chow groups, after every product; hence they
are inverse motivic maps. See \citet[Theorem 2.8(3), p.~1472,
and Corollary 3.10, p.~1484]{Jiang2023}.
\end{proof}

\begin{lemma}\label{lem:chow-descent}
Let $k'/k$ be a finite extension of finite fields. Suppose that
$\mathsf h(X_{k'})$ is a summand of a finite direct sum of motives
$\mathsf h(A_j)\otimes\mathbb L^{r_j}$, where the $A_j$ are abelian
varieties over $k'$ and $r_j\ge0$. Then $\mathsf h(X)$ is a summand
of a finite direct sum of abelian, Artin, and Lefschetz motives over
$k$. This holds for Chow retractions and also for correspondences
whose retraction identities hold in every $\ell$-adic realization.
\end{lemma}

\begin{proof}
Regard a $k'$-variety as a $k$-variety and a correspondence over $k'$
as a correspondence over $k$, by extension by zero along the open
and closed inclusion $Y\times_{k'}Z\subset Y\times_k Z$.
This defines induction of motives and preserves composition: after
base change to $k'$, it is the direct sum of the conjugate
correspondences, with zero off the diagonal.

For $A/k'$, put $B=\operatorname{Res}_{k'/k}A$. This is an abelian
variety: the Weil restriction exists for quasiprojective schemes
\citep[\S4.6, Proposition 4.6.3, pp.~110--111]{Poonen2017}, and after
the separable extension $k'/k$ it is the product of the conjugates
of $A$. Projection onto the identity factor and insertion using the
origins in the other factors split $\mathsf h(A)$ from
$\mathsf h(B_{k'})$. Induction therefore splits
$\operatorname{Ind}_{k'/k}\mathsf h(A)$ from
$\mathsf h(B)\otimes\mathsf h(\operatorname{Spec}k')$.
Finally the finite map $X_{k'}\to X$ gives a retraction after dividing
its pushforward by $[k':k]$. Composing these splittings proves the
claim. All identities used are correspondence identities, so the
same argument applies to realization-level retractions.
\end{proof}

\subsection{Kummer surfaces and separated cubic fourfolds}

The following two families admit the required correspondences
directly from their geometry.

\begin{proposition}\label{prop:kummer}
Let $A$ be an abelian surface over a finite field of characteristic
different from two, and let $K=\operatorname{Km}(A)$ be its Kummer
surface. There is a Chow retraction
\[
 \mathsf h(K)\ \mid\
 \mathsf h(A)\oplus\mathsf h(A[2])\otimes\mathbb L.
\]
If $A=E_1\times E_2$, all powers of $K$ satisfy the Tate conjecture
in every codimension over every finite extension. Its formal Brauer
height is one if both elliptic curves are ordinary, and two if exactly
one is ordinary and the other supersingular.
\end{proposition}

\begin{proof}
Blow up the finite etale subscheme $A[2]$. The quotient of
$\widetilde A=\operatorname{Bl}_{A[2]}A$ by the lifted involution
is the smooth Kummer surface, and the quotient map has degree two.
This is the usual Kummer construction
\citep[Chapter 1, Example 1.3(iii), author draft pp.~8--9]{Huybrechts2016}.
Lemma~\ref{lem:retractions} gives the asserted Chow retraction.

For the product case, every power of the right-hand motive is a
summand of sums of motives of products of elliptic curves, with
Artin and Lefschetz factors. Tate classes on products of elliptic
curves are generated by divisor classes; see
\citet[Appendix A.7, first example, pp.~158--159]{Milne2001}.
An Artin factor becomes a sum of units after a finite extension, and
averaging cycles over the finite Galois group descends the Tate
statement. Taking tensor powers of the retraction proves it for $K$.

For the height calculation, the degree-two pullback identifies
$H^2(K)$ with $H^2(A)\oplus H^0(A[2])(-1)$; it is injective and
both sides have dimension $22$. The same holds for crystalline
cohomology. The Artin summand contributes sixteen slopes equal to
one, whether or not the two-torsion points are rational. The slopes
of an ordinary elliptic curve are $0,1$,
and those of a supersingular elliptic curve are $1/2,1/2$.
The Kunneth formula therefore gives
\[
 \begin{array}{c|c}
 (E_1,E_2)&\text{slopes of }H^2(K)\\ \hline
 \text{ordinary, ordinary}&0,\ 1^{20},\ 2\\
 \text{ordinary, supersingular}&(1/2)^2,\ 1^{18},\ (3/2)^2.
 \end{array}
\]
Here a superscript denotes multiplicity. The finite-height K3
Newton polygon has slopes $1-1/h$ and $1+1/h$, each with
multiplicity $h$, and $1$ with multiplicity $22-2h$;
see \citet[Chapter 18, \S3.2--3.3, author draft pp.~399--401]{Huybrechts2016}.
Thus the displayed heights are respectively one and two.
\end{proof}

\begin{proposition}\label{prop:cubic-join}
Let $k$ be a finite field of characteristic different from three.
Let $f$ and $g$ be ternary cubic forms defining smooth plane cubic
curves $E_f,E_g$, and set
\[
 X:\quad f(x_0,x_1,x_2)+g(y_0,y_1,y_2)=0\ \subset\mathbf P^5.
\]
Then $X$ is smooth. After a finite extension, its Chow motive is a
summand of
\begin{equation}\label{eq:cubic-target}
 (\mathbf1\oplus\mathbb L^{\oplus7}\oplus\mathbb L^2)^{\otimes2}
       \ \oplus\ \mathsf h(E_f\times E_g)\otimes\mathbb L.
\end{equation}
Over $k$, it is a summand of abelian, Artin, and Lefschetz motives.
Every power of $X$ satisfies the Tate conjecture in every codimension
over every finite extension.
\end{proposition}

\begin{proof}
Euler's identity and $3\ne0$ show that the gradient of $f$ has
no zero at a nonzero vector, and likewise for $g$. This proves
smoothness of $X$ and of the cubic surfaces
\[
 S_f:\ f(x)+u^3=0,\qquad S_g:\ g(y)+v^3=0.
\]
On their product consider
\[
 ((x:u),(y:v))\dashrightarrow(xv:-uy).
\]
Its image lies in $X$, since
$f(xv)+g(-uy)=v^3(-u^3)-u^3(-v^3)=0$.
The base locus is $E_f\times E_g$, and locally its base ideal is
$(u,v)$: some coordinate of $x$ and some coordinate of $y$ are units.
Consequently its blowup
\[
 V=\operatorname{Bl}_{E_f\times E_g}(S_f\times S_g)
\]
resolves the map. Above a point $(a:b)\in X$ with $f(a)\ne0$,
the preimages are
\[
 ((a:u),(-b/u:1)),\qquad u^3=-f(a).
\]
There are three geometric preimages and the equation is separable.
Thus $V\to X$ is generically finite of degree three. This is the
degree-three instance of the construction of
\citet[Remark 1.10, p.~102]{ShiodaKatsura1979}.

Over a finite extension each cubic surface is a plane blown up at
six points, giving
$\mathsf h(S_f)\simeq\mathsf h(S_g)\simeq
\mathbf1\oplus\mathbb L^{\oplus7}\oplus\mathbb L^2$.
The classification used here is
\citet[Theorem 9.4.4, p.~276]{Poonen2017}.
Lemma~\ref{lem:retractions} now gives \eqref{eq:cubic-target},
with retraction $\frac13\phi_*$. Over a further finite extension
the plane cubics have origins and are elliptic curves.
Lemma~\ref{lem:chow-descent} descends the retraction to $k$.
The Tate assertion follows from the elliptic-product theorem exactly
as in Proposition~\ref{prop:kummer}.
\end{proof}

These are explicit constructions for the stated families. In
particular Proposition~\ref{prop:cubic-join} does not cover an
arbitrary smooth cubic fourfold.

\subsection{The exact mixed correspondence condition}

Fix a finite field $k=\F_q$. A \emph{target with the full Tate
property} will mean a motive in the tensor category generated by
a specified finite collection of abelian varieties, finite etale
schemes, and Lefschetz motives, such that every product of the
abelian generators satisfies the Tate conjecture in every
codimension after every finite extension. This hypothesis is not
asserted for arbitrary collections of abelian varieties.
In such a category, cycle classes span all Frobenius-equivariant
maps between realizations: apply Kunneth and Poincare duality to
the product presenting the two motives, and then apply their
algebraic projectors. Duals and Tate twists give the same statement
for maps to $\mathbb L^a$.

\begin{theorem}\label{thm:mixed-condition}
Let $X/k$ be a finite-height K3 surface or a smooth cubic fourfold,
and let $d=\dim X$. Suppose that a target with the full Tate property
has a finite presentation
\[
 M=(T,e),\qquad T=\bigoplus_{\nu=1}^s\mathsf h(Y_\nu),
\]
where $e$ is a rational algebraic projector in all realizations,
the $Y_\nu$ are products of the target generators and projective
spaces, and $M$ has the same Frobenius characteristic polynomial
as $X$ in every cohomological degree. The following are equivalent.
\begin{enumerate}
\item There are rational algebraic correspondences
$f:\mathsf h(X)\to M$ and $g:M\to\mathsf h(X)$ that are inverse
in every $\ell$-adic realization, $\ell\ne p$.
\item For one prime $\ell_0\ne p$, the projected map
\begin{equation}\label{eq:mixed-map}
 \bigoplus_\nu\operatorname{CH}^{d}(X\times Y_\nu)_{\Q}
       \otimes\Q_{\ell_0}
 \longrightarrow
 \operatorname{Hom}_{F}^{\mathrm{gr}}
       (H^*_{\ell_0}(X),H^*_{\ell_0}(M)),\quad z\longmapsto e_*z_*,
\end{equation}
is surjective.
\item For one prime $\ell_0\ne p$, the Tate conjecture in
codimension $d$ holds for each of the finitely many products
$X\times Y_\nu$.
\end{enumerate}
Here the maps on the right of \eqref{eq:mixed-map} preserve degree
and commute with Frobenius. If these conditions hold at one prime,
they hold at every prime different from $p$.
\end{theorem}

\begin{proof}
Frobenius is semisimple on the cohomology in question. For K3
surfaces, see the semisimplicity input in
\citet[Lemma 10.4 and Remark 10.5, p.~54]{IIK2021}; for cubic
fourfolds, see \citet[introduction, author preprint p.~1]{Zarhin2004}.
Abelian varieties and the other target generators also have
semisimple Frobenius. Thus equality of the degreewise characteristic
polynomials supplies an isomorphism at $\ell_0$.

Assume (2). Lift such an isomorphism to an algebraic correspondence
$f_0$ with $\Q_{\ell_0}$ coefficients. Its inverse is algebraic
with these coefficients as well, for the following reason, which
does not presume any mixed Tate statement in the reverse direction.
Transport the Poincare duality form of $X$ by $f_0$ to a
Frobenius-equivariant pairing
$H^*(M)\otimes H^*(M)\to\Q_{\ell_0}(-d)$.
The full Tate property of the target lifts this pairing to a
correspondence. Compose it with the dual of $f_0$ and the inverse
of the canonical motivic duality of $X$. This gives a correspondence
$g_0:M\to\mathsf h(X)$ realizing $f_0^{-1}$.

Write $f_0,g_0$ as finite $\Q_{\ell_0}$-linear combinations of
rational cycles. Nonvanishing of the determinant of the composition
is an open condition on their coefficient tuples. Since $\Q$ is
dense in $\Q_{\ell_0}$, there are rational combinations $f,g$ for
which $u=gf$ is invertible at $\ell_0$.

Both types of $X$ have zero odd cohomology. For each $r\ge1$,
the correspondence trace formula therefore reads
\begin{equation}\label{eq:trace-certificate}
 \operatorname{Tr}(u_*^r\mid H^*_\ell(X))
           =\deg(\Delta_X\cdot u^r)\in\Q.
\end{equation}
The right side is independent of $\ell$. The formula follows
directly by expanding the diagonal in dual Kunneth bases; see
\citet[Lemmas 45.9.7 and 45.10.3, Tags 0FHH and 0FHR]{Stacks}.
Newton identities give a common rational characteristic polynomial
$P_u(T)$ for $u_*$, of degree $\dim H^*(X)$. Its constant term
is nonzero, because it is nonzero at $\ell_0$. Cayley--Hamilton
expresses $u_*^{-1}$ as a rational polynomial $R(u_*)$ at every
prime. Replace $g$ by $R(u)g$. Then $gf=1$ in all realizations;
equal dimensions give $fg=1$ there as well. This proves (1).

If (1) holds, any Frobenius-equivariant graded map from $H^*(X)$
to $H^*(Y_\nu)$ is the composite of $f$ with a map from $M$ to
$\mathsf h(Y_\nu)$. The full Tate property lifts the latter map.
Kunneth and Poincare duality identify all these maps with
codimension-$d$ Tate classes on $X\times Y_\nu$. Thus (1) implies
(3) at every prime. Finally (3) implies (2) by projection with $e$.
\end{proof}

The condition can be restricted to the remaining summand.
For a cubic fourfold, let $H$ be its hyperplane class and put
\[
 p_{2i}=\tfrac13 H^{4-i}\times H^i\quad(0\le i\le4),
 \qquad p_{\mathrm{prim}}=\Delta_X-\sum_{i=0}^4p_{2i}.
\]
These are mutually orthogonal Chow projectors, since $\deg H^4=3$;
the last one realizes the $22$-dimensional primitive $H^4$.
For a K3 surface, pass to a finite extension defining a basis
$D_1,\ldots,D_\rho$ of the Neron--Severi group over $\Q$.
Choose a rational zero-cycle $o$ of degree one, let
$B=(D_i\cdot D_j)$, and put
\[
 p_0=o\times X,\quad p_4=X\times o,\quad
 p_{\mathrm{NS}}=\sum_{i,j}(B^{-1})_{ij}D_i\times D_j,\quad
 p_{\mathrm{tr}}=\Delta_X-p_0-p_4-p_{\mathrm{NS}}.
\]
The intersection form is nondegenerate, and multiplication of
product correspondences verifies all projector identities in Chow
groups. The complement is
$\mathbf1\oplus\mathbb L^{\oplus\rho}\oplus\mathbb L^2$.

If a target $N=(T,e)$ with the full Tate property has the same
realization as $p_{\mathrm{prim}}$ or $p_{\mathrm{tr}}$, respectively,
the exact remaining condition is the surjectivity of
\[
 z\longmapsto e_*z_*p_*:
 \bigoplus_\nu\operatorname{CH}^{d}(X\times Y_\nu)_{\Q}
            \otimes\Q_{\ell_0}
 \longrightarrow
 \operatorname{Hom}_F(H^*_{\ell_0}(X,p),H^*_{\ell_0}(N)).
\]
Indeed the known complement already has algebraic comparison maps.
Adjoining these maps reduces the assertion to
Theorem~\ref{thm:mixed-condition}. Conversely an isomorphism of the
full motives transports $p$ to an algebraic projector on the target;
the full Tate property supplies the comparison with the specified
$N$. This also explains why the choice of a presentation is harmless.
Here target existence is a hypothesis; no general CM realization
theorem is being assumed without its hypotheses.

\subsection{A target for a weight-two slope pattern}

The ordinary case permits a target construction using only
Honda--Tate theory and a published all-powers Tate theorem. We give
the arithmetic details, including the possible noncommutative
endomorphism algebra.

\begin{proposition}\label{prop:ordinary-packet}
Let $f(T)\in\Z[T]$ be a nonconstant monic polynomial whose roots
are algebraic integers of complex absolute value $q$, where
$q=p^a$. With valuations normalized by $v(q)=1$, suppose that its
roots have slopes $0$ and $2$ each once and slope $1$ otherwise.
Assume that no root is $q$ times a root of unity. Then:
\begin{enumerate}
\item $f$ is irreducible and separable. For a root $\alpha$,
$K=\Q(\alpha)$ is a CM field and $\Q(\alpha^r)=K$ for every $r>0$.
\item There is a geometrically simple abelian variety $B/\F_{q^2}$
with Frobenius $\alpha$. Its endomorphism division algebra $D$ has
center $K$ and degree $d\in\{1,2\}$ over $K$.
\item Every power of $B$ satisfies the Tate conjecture in every
codimension after every finite extension. There is a rational
algebraic homological summand $N$ of $\mathsf h^1(B)^{\otimes2}$
whose Frobenius roots, each once, are
\[
                  \{\sigma(\alpha)^2:\sigma:K\hookrightarrow\overline\Q\}.
\]
All the projectors defining $N$ are the same rational
correspondences for every $\ell\ne p$.
\end{enumerate}
\end{proposition}

\begin{proof}
Choose a root of slope zero. Its irreducible factor also contains
its complex conjugate $q^2/\alpha$, of slope two. That factor has
multiplicity one. Any other factor has all its roots of slope one
at all primes above $p$. For one of its roots $\beta$, the number
$\beta/q$ is an algebraic integer: this is clear away from $p$,
and follows from the slope assertion above $p$. Every complex
conjugate has absolute value one, so Kronecker's theorem makes
$\beta/q$ a root of unity, contrary to the assumption. Thus $f$
is irreducible. Conjugation acts on $K$ by
$\alpha\mapsto q^2/\alpha$, commutes with every embedding, and
has a totally real fixed field. No embedding of $K$ is real,
since a real conjugate would be $q$ or $-q$. Hence $K$ is CM.

For $L=\Q(\alpha^r)$, use the slope-zero embedding of $K$ into
$\overline\Q_p$. Every $L$-embedding of $K$ sends $\alpha$ to a
root with the same $r$th power, hence to another slope-zero root.
There is only one. Separability gives $[K:L]=1$.
The prime of $K$ where $\alpha$ has slope zero has local degree
one: the number of embeddings inducing that prime equals its local
degree and each would give a slope-zero root. Its conjugate also
has local degree one.

Now $\alpha$ is a Weil $q^2$-number. Its slopes relative to $q^2$
are $0,1/2,1$. The Honda--Tate local invariants have denominators
at most two, so the division-algebra degree is $d=1$ or $2$.
The equality $\Q(\alpha^r)=K$ and constancy of these invariants
after every extension prove geometric simplicity and show that
all geometric endomorphisms are already present.

For completeness, the all-powers Tate statement can be seen from
the same unique slope-zero prime. A conjugate of $\alpha/q$ has
valuations supported at just the slope-zero and slope-two
embeddings, with values $-1,+1$. If
\[
        \prod_\sigma (\sigma(\alpha)/q)^{m_\sigma}
\]
is a root of unity, applying all conjugate valuations forces
$m_\sigma=m_{\bar\sigma}$ for each conjugate pair. Thus every
Tate eigenmonomial in the exterior algebra on any number of copies
of $H^1(B)$ is a product of degree-two Tate eigenmonomials.
The divisor Tate theorem makes those factors algebraic. Scalar
extension to a splitting field and descent of the span give the
assertion over $\Q_\ell$. This is precisely the second example of
\citet[Appendix A.7, p.~159]{Milne2001}, which also covers $d=2$;
it is not restricted to the commutative-endomorphism case.

It remains to remove the multiplicity $d$ in $H^1(B)$.
An algebraic projector onto $\mathsf h^1(B)$ is obtained by
polynomial interpolation in $[m]^*$, whose eigenvalues on
$H^i(B)$ are $m^i$, for any integer $m>1$.
The separable algebra $K\otimes_\Q K$ has an idempotent
$e_\Delta$ for its multiplication factor $K$; over a splitting
field it keeps exactly the pairs of equal embeddings.
Inside the algebra of endomorphisms of
$\mathsf h^1(B)^{\otimes2}$, its block is
\[
 e_\Delta(D\otimes_\Q D)\simeq D\otimes_K D.
\]
Every local invariant of $D$ has order at most two. The last
algebra therefore has trivial Brauer class and is
$\operatorname{Mat}_{d^2}(K)$. Choose a primitive idempotent $j$
in it. On the equal-embedding block the realization has dimension
$d^2$ at each embedding of $K$, and this matrix algebra acts by
its standard module. The image of $j$ consequently has dimension
one at each embedding and Frobenius $\sigma(\alpha)^2$ there.
The corresponding rational algebraic projector defines $N$.
An opposite-algebra convention for $\mathsf h^1$ changes none of
this argument, since the Brauer class has order at most two.
The asserted idempotences hold modulo homological equivalence;
no lifting to Chow idempotents is asserted here.
\end{proof}

\subsection{Ordinary K3 surfaces and ordinary cubic fourfolds}

We apply the slope-pattern construction after removing the
algebraic part of the middle cohomology.

\begin{proposition}\label{prop:ordinary-target}
Let $X/\F_q$ be an ordinary K3 surface or an ordinary smooth cubic
fourfold, in any characteristic. After a finite extension, there is
an explicitly specified homological motive $M$ with the full Tate
property whose graded Frobenius realization equals that of $X$.
It is built from powers of a single abelian variety, Lefschetz
motives, and direct summands. For this target the remaining
assertion that $\mathsf h(X)\simeq M$ by rational algebraic
correspondences is exactly the mixed condition of
Theorem~\ref{thm:mixed-condition}.
\end{proposition}

\begin{proof}
The divisor Tate theorem for K3 surfaces holds also in
characteristics two and three; it follows, in particular, from
\citet[Theorem 1.5, p.~3, and Lemma 10.3, pp.~53--54]{IIK2021}.
For cubic fourfolds, codimension-two Tate in odd characteristic is
\citet[Theorem 5.14(3), author preprint p.~27]{MadapusiPera2015}.
In characteristic two it is
\citet[Theorem 4.10, author-hosted journal text p.~20]{AuelEtAl2025}.
Together with the semisimplicity references in
Theorem~\ref{thm:mixed-condition}, these statements identify the
entire root-of-unity part of the middle cohomology with algebraic
classes after a finite extension. Enlarge the field so all those
classes are defined, and rename the field $\F_q$.

The middle Frobenius polynomial is rational. Its roots of the
form $q^r\zeta$, with $\zeta$ a root of unity and $r=1$ or $2$,
form a Galois-stable multiset. Removing their factor therefore
leaves a monic integral polynomial. Semisimplicity and Poincare
duality make the pairing on the algebraic eigenspace nondegenerate.
Choose rational algebraic cycles $C_1,\ldots,C_t$ whose classes
form a basis of this eigenspace at one prime, and put
$G=(C_i\cdot C_j)$. The rational matrix $G$ is nonsingular, so
the classes remain independent at every prime; the characteristic
polynomial shows that they span the same-dimensional eigenspace
there. Their Chow projector is
$\sum_{i,j}(G^{-1})_{ij}C_i\times C_j$.
Its complement realizes exactly the remaining polynomial.

For an ordinary K3 surface the slopes of $H^2$ are
$0,1^{20},2$. Removing the algebraic factor removes only slope-one
roots. The remaining polynomial meets
Proposition~\ref{prop:ordinary-packet}. After the quadratic
extension $\F_{q^2}$, its nonalgebraic realization is therefore
that of the motive $N$ constructed there. Adjoin
$\mathbf1\oplus\mathbb L^{\oplus\rho}\oplus\mathbb L^2$.

For an ordinary cubic fourfold the middle slopes are
$1,2^{21},3$; see \citet[\S4, discussion after Remark 4.1,
author preprint p.~6]{Zarhin2004}.
Write a remaining root as $\gamma$ and put $\alpha=\gamma/q$.
This is an algebraic integer: outside $p$, division by $q$ is
harmless; at every prime above $p$, all conjugates of $\gamma$
have valuation at least $v(q)$. The polynomial of these
$\alpha$'s is consequently monic integral, of weight two, with
slopes $0,1,2$ as in Proposition~\ref{prop:ordinary-packet}, and
without roots $q\zeta$. Over $\F_{q^2}$, the motive
$N\otimes\mathbb L$ has roots $q^2\alpha^2=\gamma^2$.
Adjoin the algebraic middle summand and the known outer-degree
Lefschetz summands to obtain $M$.

All the required Lefschetz motives occur in powers of the same
$B$: an ample divisor and a suitably normalized complementary
power split $\mathbb L$ from $\mathsf h(B)$. Hence $M$ has a
finite untwisted presentation using powers of $B$, and
Proposition~\ref{prop:ordinary-packet} supplies its full Tate
property. The exact remaining condition is now
Theorem~\ref{thm:mixed-condition}. If the comparison
correspondences are constructed over an extension, induction,
Weil restriction, and trace as in Lemma~\ref{lem:chow-descent}
give a comparison with an abelian--Artin target over the original
field. The conjugates of $B$ have the same Frobenius polynomial
and are isogenous, so this descent does not introduce an
unproved mixed Tate assertion among unrelated abelian varieties.
\end{proof}

The target construction and the mixed comparison are different
conclusions. The proposition proves the former. It does not
construct the missing mixed cycles for an arbitrary ordinary K3
surface or ordinary cubic fourfold. The known Tate theorem for
$X\times X$ in the K3 case
\citep[Theorem 1.5]{IIK2021}, and for all self-powers in the
ordinary cubic case \citep[Corollary 5.2, author preprint p.~8]{Zarhin2004},
do not by themselves imply Tate for $X\times Y_\nu$ with the
specified abelian targets. The retractions of
Propositions~\ref{prop:kummer} and \ref{prop:cubic-join} supply
such cycles for their concrete families; the general mixed
condition remains explicit.


\section{What remains open}

For a general simple abelian variety, the least exponent must allow
Frobenius fields larger than the eventual geometric center and must
preserve the multiplicity of the geometric simple factor. Even in
the imaginary quadratic setting, Proposition~\ref{prop:central}
does not assert that every minimum is attained by a central
Frobenius. The remaining problem is to characterize the noncentral
elements of $D$ that can represent a compatible Weil number over a
smaller field. Theorem~\ref{thm:main} resolves that obstruction for
one specific geometric isogeny class. Proposition
\ref{prop:finite-test} computes the minimum from a given Frobenius
polynomial. Proposition~\ref{prop:central-cm} computes every
central minimum and, when the geometric endomorphism algebra is
commutative, the unrestricted minimum. Its general sufficient
bound includes the residue degrees and is not proved sharp. A
sharp unrestricted bound using only class numbers and Brauer
indices is not established here.

For general K3 surfaces and cubic fourfolds, the remaining issue is
the construction of the mixed correspondences in
Theorem~\ref{thm:mixed-condition}. Proposition~\ref{prop:ordinary-target}
constructs a target in the ordinary case and proves its requisite
Tate property. It does not settle that mixed condition. The explicit
Kummer and separated-cubic families in Section~\ref{sec:companion}
are cases where geometric maps already provide the correspondences.

\paragraph{Method and availability.}
GPT-6 Astra (ChatGPT Pro mode) produced the mathematics and text.
The concrete threefold and its failure of descent to $\F_8$ were
re-proved from the statement before comparison with earlier
arguments, and the full argument was subsequently checked by the
same model. These are same-model checks, not independent human
verification. The exact minimum exponent was also separately
re-proved from its statement by the same model, and the full
argument received a further targeted same-model review. The
general CM central-minimum formula, the Kummer and separated-cubic
retractions, the mixed correspondence
criterion, and the weight-two slope-pattern proposition were each
separately re-proved from their statements by the same model
before comparison with the written proofs. The companion results
also received targeted proof reviews, including their descent and
ordinary-target applications. Their complete dependency arguments
and the primary references specified below were checked.
No human has checked this paper. The text is dedicated
to the public domain under
\href{https://creativecommons.org/publicdomain/zero/1.0/}{CC0 1.0}.
Sources, checks, and editable material are associated with the
\href{\workbenchurl}{algebraic-groups-and-representations workbench}; the
permanent address for its records is
\href{\recordurl}{mathematics-workbenches}.

\clearpage
\begin{thebibliography}{99}

\bibitem[Auel et al.(2025)]{AuelEtAl2025}
Auel, A., Kulkarni, A., Petok, J., \& Weinbaum, J. (2025).
A census of cubic fourfolds over $\F_2$.
\emph{Mathematics of Computation, 94}, 2089--2112.
\url{https://doi.org/10.1090/mcom/4010}
Theorem 4.10, author-hosted journal text p.~20 (the theorem numbering
differs from arXiv:2306.09908v1).
\url{https://jonahweinbaum.github.io/material/cubic_fourfolds.pdf}

\bibitem[Brooks et al.(2017)]{BrooksEtAl2017}
Brooks, E. H., Jetchev, D., \& Wesolowski, B. (2017).
Isogeny graphs of ordinary abelian varieties.
\emph{Research in Number Theory, 3}, Article 28.
Lemma 5.5 and \S5.3: unit norms and the Shimura class group.
\url{https://doi.org/10.1007/s40993-017-0087-5}

\bibitem[Dupuy et al.(2020)]{DKRV2020}
Dupuy, T., Kedlaya, K. S., Roe, D., \& Vincent, C. (2020).
\emph{Isogeny classes of abelian varieties over finite fields in the
LMFDB} (arXiv:2003.05380v2, September 23).
\S3.1, pp.~12--13; \S3.4, pp.~14--15, especially
Propositions 3.2 and 3.5.
\url{https://arxiv.org/abs/2003.05380v2}

\bibitem[Grothendieck(n.d.)]{GrothendieckND}
Grothendieck, A. (n.d.).
[Notes on motives over finite fields] [Unpublished manuscript].
Archives Grothendieck, Universit\'e de Montpellier, cote 15,
scans 94--106, especially scans 100, 102, 104, and 106
(manuscript footers 99, 101, 103, and 105).
\url{https://grothendieck.umontpellier.fr/}

\bibitem[Huybrechts(2016)]{Huybrechts2016}
Huybrechts, D. (2016).
\emph{Lectures on K3 surfaces}. Cambridge University Press.
Chapter 1, Example 1.3(iii), and Chapter 18, \S3.2--3.3;
author draft pp.~8--9 and 399--401, respectively.
\url{https://www.math.uni-bonn.de/people/huybrech/K3Global.pdf}

\bibitem[Ito et al.(2021)]{IIK2021}
Ito, K., Ito, T., \& Koshikawa, T. (2021).
CM liftings of K3 surfaces over finite fields and their applications
to the Tate conjecture.
\emph{Forum of Mathematics, Sigma, 9}, e29, 1--70.
\url{https://doi.org/10.1017/fms.2021.24}
Theorem 1.5, p.~3; Lemmas 10.3--10.4 and Remark 10.5, pp.~53--54.

\bibitem[Jiang(2023)]{Jiang2023}
Jiang, Q. (2023).
On the Chow theory of projectivizations.
\emph{Journal of the Institute of Mathematics of Jussieu, 22}(3),
1465--1508.
\url{https://doi.org/10.1017/S1474748021000451}
Theorem 2.8(3), p.~1472; Corollary 3.10, p.~1484.

\bibitem[Madapusi Pera(2015)]{MadapusiPera2015}
Madapusi Pera, K. (2015).
The Tate conjecture for K3 surfaces in odd characteristic.
\emph{Inventiones Mathematicae, 201}, 625--668.
\url{https://doi.org/10.1007/s00222-014-0557-5}
Theorem 5.14(3), author preprint p.~27.
\url{https://arxiv.org/abs/1301.6326}

\bibitem[Milne(2001)]{Milne2001}
Milne, J. S. (2001).
The Tate conjecture for certain abelian varieties over finite fields.
\emph{Acta Arithmetica, 100}(2), 135--166.
\url{https://doi.org/10.4064/aa100-2-3}
Appendix A.7, first and second examples, pp.~158--159.
The author's erratum concerns Theorem B.1, not these examples.
\url{https://www.jmilne.org/math/articles/2001a.html}

\bibitem[Milne(2020)]{Milne2020}
Milne, J. S. (2020).
\emph{Algebraic number theory} (Version 3.08, July 19).
Theorem 4.3, p.~70.
\url{https://www.jmilne.org/math/CourseNotes/ANT.pdf}

\bibitem[Poonen(2017)]{Poonen2017}
Poonen, B. (2017).
\emph{Rational points on varieties}. American Mathematical Society.
\S4.6, Proposition 4.6.3, pp.~110--111;
Theorem 9.4.4, p.~276.
\url{https://math.mit.edu/~poonen/papers/Qpoints.pdf}

\bibitem[Shioda and Katsura(1979)]{ShiodaKatsura1979}
Shioda, T., \& Katsura, T. (1979).
On Fermat varieties.
\emph{Tohoku Mathematical Journal, 31}(1), 97--115.
Remark 1.10, p.~102.
\url{https://www.jstage.jst.go.jp/article/tmj1949/31/1/31_1_97/_pdf}

\bibitem[The Stacks Project Authors(2026)]{Stacks}
The Stacks Project Authors. (2026).
\emph{The Stacks Project}. Lemma 11.7.4, Tag 074W;
Lemmas 45.9.7 and 45.10.3, Tags 0FHH and 0FHR.
\url{https://stacks.math.columbia.edu/tag/074W}
\url{https://stacks.math.columbia.edu/tag/0FHH}
\url{https://stacks.math.columbia.edu/tag/0FHR}

\bibitem[Tate(1971)]{Tate1971}
Tate, J. (1971).
Classes d'isog\'enie des vari\'et\'es ab\'eliennes sur un corps fini
(d'apr\`es T.~Honda).
In \emph{S\'eminaire Bourbaki, 1968/69}, expos\'e 352
(pp.~95--110). Springer.
Theorem 1(i)--(ii) and Remark 2, p.~96.
\url{https://numdam.org/item/SB_1968-1969__11__95_0/}

\bibitem[Thakur(2020)]{Thakur2020}
Thakur, S. (2020).
\emph{Isogenies of certain abelian varieties over finite fields with
$p$-ranks zero} (arXiv:2001.00210v1, January 1).
Definition 2.1, p.~3; Proposition 2.4, pp.~4--5;
Definition 2.3 and Proposition 2.9, p.~6.
\url{https://arxiv.org/abs/2001.00210v1}

\bibitem[Zarhin(2004)]{Zarhin2004}
Zarhin, Y. G. (2004).
The Tate conjecture for powers of ordinary cubic fourfolds over
finite fields.
\emph{Journal of Number Theory, 108}, 44--59.
\url{https://doi.org/10.1016/j.jnt.2004.04.006}
Introduction, \S4 after Remark 4.1, and Corollary 5.2;
author preprint pp.~1, 6, and 8, respectively.
\url{https://arxiv.org/abs/math/0304014}

\end{thebibliography}
\end{document}
